\section{Limitations}
The benchmark covers six small parabolic equations that share two fields. It leaves out turbulence, shocks, moving geometry, mixed physical units, and equations that need other kinds of boundary data. This keeps the comparisons easy to read and to reproduce, but it limits how far the results generalize. They describe agreement under this setup, not a universal ranking of PDEs.

The reference has a strong structure of its own. The initial fields vanish near the boundary and the time horizon is short, so changing the BC moves a solution much less than changing the equation (\cref{fig:configurations}). Rankings within an IC family are therefore mostly rankings of equations, and the split into same-BC and different-BC pairs should be read with this in mind. Other horizons, field scales, or IC distributions would change the rankings.

The embedding results use 10 of the 100 samples per problem, chosen in advance because only a two-core CPU was available. The reference rankings on these ten samples agree closely with those on all 100, but the bootstrap intervals rest on only ten samples. All three models ran in bfloat16 on a CPU. The checks in \cref{app:embedding-protocol,app:extra} suggest that this changes little, but a GPU run at each model's reference precision would be a useful replication.

Five snapshots are a compressed view of a solution. Measuring the reference on the same snapshots avoids giving the image models credit for information they never saw, and the 100-frame reference shows the cost of this compression. The image also fixes one way of drawing the data: a model can only see differences that survive the colormap, the $[0,0.8]$ range, and its own resizing.

Finally, the models were pretrained on large collections of text and images that may include PDEs. Their training data are not documented in enough detail to rule this out. New random ICs, controlled coefficients, and equivalent rewrites reduce the chance that a model has seen these exact problems, but they cannot prove that an equation was absent from pretraining.

\section{Conclusion}
We asked whether general-purpose multimodal embeddings rank PDE problems by how their solutions behave. To answer it, we simulated 5,400 trajectories with one validated solver, defined a reference distance whose rankings are reliable and partly proved, and queried three embedding models with descriptions, images, or both. On ten samples per problem, chosen in advance, image embeddings follow the reference moderately well (Spearman \spImageMin--\spImageMax{}) but do no better than raw pixels. Description embeddings follow it weakly (\spTextMin--\spTextMax{}), close to character $n$-grams, and for two of the models an equivalent rewrite moves a description further than any change of problem. Joint inputs stay close to the text. These models recognize solutions that \emph{look} alike much better than equations that \emph{behave} alike. Natural next steps are the full 100-sample run on a GPU, more coefficients and time horizons, and encoders trained with physical supervision, all evaluated under the same fixed protocol.
